\subsection{对应的性质 III：非分歧主系列}

\label{s:unramps}
\ssszh 有了 $p$-adic 群的万有表示，可以如下造出更多表示。如同表示论的其他部分，有抛物诱导函子，可以从较小群（例如 Cartan 群 $T$）的表示出发构造 $G$ 的表示。把非分歧表示代入这一构造，就得到一族重要的模，即非分歧主系列。它们的单子商恰是那些关于紧开 Iwahori 子群 $I$ 有不平凡不变量的单 $G(\Fqt)$-模；这里把 $I$ 理解为 $G(\Fq[[t]])$ 在投影 $G(\Fq[[t]]) \rightarrow G(\Fq)$ 下对 $G(\Fq)$ 的 Borel 子群的原像。\footnote{这条定理常被归于 Borel 与 Matsumoto，我们也听到有人把它归于 Bernstein 或 Casselman；参见 \cite{cass80spherical} 的命题 2.4 及其中的参考文献。}也就是说，把抛物诱导作用于非分歧表示，得到的是下一类最简单的 $p$-adic 群表示。

用谱侧的语言说，$G(\Fqt)$ 的这类不可约表示对应于野分歧、幂幺单值（unipotent monodromy）的 Galois 表示。为解释这一点，从序列 \eqref{e:galunramified} 回忆惯性子群
%
\[
   I \simeq \on{Gal}( \overline{\mathbb{F}}_q(\!(t)\!)^s / \overline{\mathbb{F}}_q(\!(t)\!)),
\]
%
其中 $\overline{\mathbb{F}}_q$ 表示 $\Fq$ 的代数闭包。分歧滤过的下一个商把它分裂为
%
\[
	    1 \rightarrow I_{\on{wild}} \rightarrow I \rightarrow I_{\on{tame}} \rightarrow 1.
\]
%
对野惯性 $I_{\on{tame}}$，基本想法是：它的行为与特征零的对应物一样，即与 $k$ 上形式穿孔圆盘的覆叠一样。这些覆叠的分类方式与拓扑中圆盘的覆叠空间完全一样，即对正整数 $n$，添加 $t$ 的一个 $n$ 次根。在正特征情形，它们可分当且仅当 $n$ 与 $q$ 的特征互素；而这样的覆叠穷尽了所有次数与 $q$ 互素的 Galois 扩张。显式地，可以取同构
%
\[
   I_{\on{tame}} \simeq \underset{\ell} \Pi \hspace{.7mm} \mathbb{Z}_\ell,
\]
%
其中 $\ell$ 取遍不等于 $\Fq$ 的特征的素数。Galois 表示称作野分歧的，如果它经过对 $I_{\on{wild}}$ 的商分解；进而称作幂幺单值的，如果 $I_{\on{tame}}$ 的任意元素（或等价地，某个拓扑生成元）都被送到 $\cG$ 的幂幺元素。



\ssszh 为几何讨论作定位，我们稍为详细地回顾一下 $p$-adic 表示的相应抛物诱导如何进行。给定 Borel 子群 $B$，其 Levi 商是抽象 Cartan 群 $T$，于是在有理点上得到对应
%
\[
	    T(\Fqt) \leftarrow B(\Fqt) \rightarrow G(\Fqt),
\]
%
其中向左的箭头是满射，核为幂幺根 $N(\Fqt)$ 的有理点；向右的箭头是闭嵌入。由此得到抛物诱导函子
%
\begin{equation} \label{e:parindpadic}
	    T(\Fqt)\mod \xrightarrow{\on{Res}} B(\Fqt)\mod \xrightarrow{\on{Ind}}
	    G(\Fqt)\mod.
\end{equation}
%
具体地说，把 $T(\Fqt)$ 的光滑表示 $\pi$ 通过让幂幺根平凡作用而膨胀成 $B(\Fqt)$ 的表示，再诱导到 $G(\Fqt)$。这一诱导与通常一样，具体地由旗流形上相应向量丛的光滑截面给出，即 $G(\Fqt) \overset{B(\Fqt)} \times \pi \rightarrow G(\Fqt)/B(\Fqt)$。



特别地，由此可知：$T$ 的任意非分歧特征标的抛物诱导都嵌入万有非分歧表示的抛物诱导，即万有主系列模
%
\[
	    \Fun\big( G(\Fqt) / T(\Fq[[t]]) \cdot N(\Fqt)\big).
\]
%

\begin{re} 对尚未具体接触过抛物诱导的读者，指出下面这一点或许有助于定位：用最高权分类约化群的不可约代数表示，正是抛物诱导的一个例子。

确实，对上文那样的代数群 $G,B,T,N$，在任意域上的代数表示范畴之间有伴随对
%
\[
	    \on{Rep}(G) \overset{\on{Res}}{\underset{\on{Ind}}{\rightleftarrows}}
	    \on{Rep}(B)
	    \overset{\on{Coinv}_N}{\underset{\on{Res}}{\rightleftarrows}}
	    \on{Rep}(T),
\]
%
其中 $\on{Res}$ 表示限制，$\on{Coinv}_N$ 表示 $N$ 的余不变量函子，$\on{Ind}$ 表示余诱导函子。把它们接起来，就得到抛物限制函子：它把 $G$-模送到其 $N$-余不变量，看作 $T$-模；还得到它的右伴随，即抛物余诱导函子：它把 $T$-模以 $N$ 的平凡作用延拓为 $B$-模，再余诱导到 $G$。

有了这些，最高权定理（可以表述为每个不可约 $G$-模都有一条 $N$-余不变量直线）特别地说明：每个不可约 $G$-模都作为某个抛物余诱导模的子模出现。此外，在特征零时，由完全可约性可知每个不可约模都与某个抛物余诱导模重合。借助「特征标的余诱导即线丛截面」这一显式模型，这就重新得到关于 $G/B$ 上线丛截面的 Borel--Weil 定理。 \end{re}


\ssszh 转向几何情形。此时自守侧类似地是半无限旗流形上的 D-模，即
		%
\[
	  \Dmod( LG / L^+ T \cdot LN),
\]
%
如前所述，它可以由 $\Dmod(\Gr_T)$ 通过一个类似于 \eqref{e:parindpadic} 的适当抛物诱导函子构造出来。

\begin{re} 关于这个范畴有一条技术性的评论。目前还不知道如何把空间 $LG/L^+T \cdot LN$ 当作代数几何的对象来处理。不过，它的 D-模范畴可以通过强制光滑下降来定义，即取 $\Dmod(LG)$ 的 $\Dmod(L^+T \cdot LN)$-不变量。已知（由 Raskin 的一个论证 \cite{raskin2ndadjoint}）这典范地等同于余不变量，也等同于 Iwahori 子群的（余）不变量，即仿射旗簇上的 D-模。这与 $p$-adic 情形中 Borel--Matsumoto 的结果平行，而且实际上是用类似方法证明的。
\end{re}



\ssszh 在谱侧，再次施加抛物诱导函子，即沿对应
 %
 \[
	    \on{LocSys}_{\check{T}} \leftarrow \on{LocSys}_{\check{B}} \rightarrow
	    \LS
 \]
 %
 对 $\QCoh(\on{pt}/\check{T})$ 作推-拉。这里指的是范畴层的拉回与推前，如注 \ref{r:pushpull2coh} 所述。\footnote{更仔细地说，前面许诺的对公式 \eqref{e:loclangnaive} 的修正将涉及从拟凝聚范畴层过渡到某些归纳凝聚范畴层，但这一修正不影响本计算。}




	  下面逐步说明如何计算。第一步是拉回，注意有笛卡儿图
	  %
	  \[
	  \xymatrix{ \on{LocSys}_{\check{N}} \ar[r] \ar[d] &
	  \on{LocSys}_{\check{B}} \ar[d] \\ \on{pt} \ar[r] &
	  \on{LocSys}_{\check{T}}.}
	  \]
	  %
	  这不过是说：一个 $\check{B}$ 局部系统，若其相伴的 $\check{T}$ 局部系统已平凡化，就等同于一个 $\check{N}$ 局部系统。更通俗地说，若转移矩阵的对角线上全是 $1$，那它就不是 $\check{B}$-丛，而是 $\check{N}$-丛！由 $\check{N}$ 的幂幺性，穿孔圆盘上 $\check{N}$-局部系统的模空间就是 $\check{\fn} / \check{N}$，同构由映射
	  %
	  \begin{equation} \label{e:rsmap}
	    \check{\fn} / \check{N} \xrightarrow{\sim}
	    \on{LocSys}_{\check{N}},
	    \quad  X \mapsto d +
	    X \cdot
	    \frac{dt}{t}.
	  \end{equation}
	  %
	 给出。由这段讨论可知，拉回为
	 %
	 \begin{equation}
	 \begin{split}
	 \QCoh(\on{LocSys}_{\check{B}}) \otimes_{\QCoh(\on{LocSys}_{\check{T}})}
	 \QCoh(\on{pt}/\check{T}) &\simeq \QCoh( \on{LocSys}_{P}  \times_{
	 \on{LocSys}_{\check{T}}} \on{pt}/\check{T}  ) \\
	 &\simeq \QCoh(\check{\fn} / \check{B} ).
	 \end{split}
	 \end{equation}
	 %


最后，推前再次给出 $\QCoh(\check{\fn} / \check{B} )$，通过投影
%
\begin{equation} \label{e:projection}
	     \check{\fn} / \check{B} \rightarrow \on{LocSys}_{\check{B}}
	     \rightarrow \LS.
\end{equation}
%
把它看作 $\LS$ 上的范畴层。把它纯粹用 $\cG$ 重写一遍会更清楚。为此，可以平凡地用旗簇的余切丛重写我们的叠，即
%
\[
	      \check{\fn} / \check{B} \simeq (\cG \overset{ \check{B}}
	      \times \check{\fn}) / \cG =:
	      \widetilde{\sN} / \cG.
\]
%
具体地说，$\widetilde{\sN}$ 是数对 $({\check{B}}', X)$ 的模空间，其中 ${\check{B}}'$ 是 Borel 子群，$X$ 是其幂幺根的元素。于是投影 \eqref{e:projection} 可以用 Springer 消解来写，即复合
%
\[
	    \widetilde{\sN}/ \cG \rightarrow \cfg / \cG \rightarrow \LS.
\]
%
这里第一个映射把数对 $(\check{B}', X)$ 送到 $X$，第二个映射的定义与 \eqref{e:rsmap} 类似（但离同构差得很远）。




\begin{re} 下文第 \ref{s:2realz} 节将会综述：$G$ 的主系列表示与 $\cG$ 的 Springer 消解上凝聚层之间的关系，最早出现在 $p$-adic 理论中。
\end{re}



\ssszh 前面两个例子，即图 \ref{f:table} 的前两行，涉及的是低分歧的对象。局部几何朗兰兹对应的这些情形虽属最基本，却已是严肃的结果。

一方面，如前所述，它们为（猜想中的）整体几何朗兰兹对应中最简单的情形提供了足够的局部输入；那里研究的是整体曲线上 $G$-丛模空间上的 D-模，可能还在有限多个点处带有抛物归约。



另一方面，它们在通常的表示论中有重要应用。一个显著的例子是对临界水平仿射李代数表示论的含义，我们将在下文第 \ref{ss:locaff} 节更详细地讨论。另一项关键应用是在过渡到模系数层之后，\footnote{严格说来，这些更自然地属于局部几何朗兰兹的 \'etale 版本或 Betti 版本，参见下文第 \ref{ss:variants} 节，因为它们有系数可变的自然类比。}对约化群的模表示论的应用；参见 \cite{finkelbergmirkovicconjecture}、\cite{mirkovic-vilonen}、\cite{acharrider}、\cite{richewilliamsontilting}、\cite{bezgaitsmiriwahotiwhittaker}、\cite{bezricheriderregcentralizer}、\cite{bezricheprinblock} 及其中的参考文献，以及综述 \cite{willimasonsurvey1}、\cite{williamsonsurvey2}。


鉴于这些进展的重要性，相当引人注目的是：它们只触及整个猜想的表层，即只等同了两侧非常小的子范畴。确实，用谱侧的话说，前述情形只涉及 $\TwoCoh(\LS)$ 中支撑在具有正则奇点与幂幺单值性的局部系统的形式邻域上的全子范畴；在整个穷尽 $\LS$ 的无限递增的轨迹序列中，这是最小的一块。



相比之下，下面要讨论的两个对应例子，涉及任意分歧的范畴环路群表示与 $\LS$ 上的范畴层。因此，这些是相当壮观的猜想，属于当下表示论这一部分中最深刻的结果之列；要把它们解决，很可能需要超出几何表示论标准工具箱的新想法。
